\documentclass[12pt]{article}
\usepackage[german]{babel}
\usepackage[parfill]{parskip}
\renewcommand{\familydefault}{\sfdefault}
\usepackage{geometry, fancyhdr}
 \geometry{
 a4paper,
 left=25mm,
 top=20mm,
 right=20mm,
 bottom=20mm
 }
\usepackage[utf8]{inputenc}
\usepackage{pdfpages}
\usepackage{amsmath,amssymb,amsfonts,amsthm,enumitem}
\newtheorem{defi}{Definition}
\newtheorem{bsp}{Beispiel}
\newtheorem{kom}{Kommentar}
\newtheorem{ass}{Eselsbrücke}
\usepackage{shadethm,multicol,float}
\usepackage{scalerel,amssymb}
\def\mcirc{\mathbin{\scalerel*{\circ}{j}}}
\def\msquare{\mathord{\scalerel*{\Box}{gX}}}
\newshadetheorem{ueb}{Übung}
\definecolor{shadethmcolor}{HTML}{EDF8FF}
\definecolor{shaderulecolor}{HTML}{45CFFF}
\setlength{\shadeboxrule}{.4pt}
\usepackage{icomma}
\newenvironment{los}{%
     \setlength\parindent{0pt}\par\medskip\textbf{Lösung.}\quad}{%
     \hfill\tiny$\blacksquare$\par\medskip}

\usepackage{kpfonts}
\usepackage{stackengine, pstricks-add}
\usepackage{calc}
\newlength\shlength
\newcommand\vv[2][0]{\setlength\shlength{#1pt}%
  \stackengine{-5.6pt}{$#2$}{\smash{$\kern\shlength%
    \stackengine{7.55pt}{$\mathchar"017E$}%
      {\rule{\widthof{$#2$}}{.57pt}\kern.4pt}{O}{r}{F}{F}{L}\kern-\shlength$}}%
      {O}{c}{F}{T}{S}}
\usepackage{pdfpages}
\allowdisplaybreaks
\usepackage{tcolorbox}
\usepackage{geometry}

\geometry{a4paper, left=25mm, right=25mm, top=25mm, bottom=25mm}

% Header und Footer konfigurieren
\pagestyle{fancy}
\fancyhf{} % alle Standardfelder leeren

% Header (zentriert)
\fancyhead[C]{,,Europaschule'' Gymnasium Gommern}

% Footer
\fancyfoot[L]{Dr. Amar Bapic}
\fancyfoot[C]{Analysis (gAn)}
\fancyfoot[R]{\thepage}

% Linien oben und unten entfernen (optional)
\renewcommand{\headrulewidth}{0.4pt}
\renewcommand{\footrulewidth}{0.4pt}

% Farben definieren
\definecolor{myblue}{RGB}{227,242,253}
\definecolor{mygreen}{RGB}{232,245,233}
\definecolor{lightorange}{RGB}{255, 213, 128}
\newcommand{\limeslinks}[2]{%
  \lim_{\substack{#1 \to #2 \\ #1 < #2}}%
}

\newcommand{\limesrechts}[2]{%
  \lim_{\substack{#1 \to #2 \\ #1 > #2}}%
}

\begin{document}

\section*{Stetigkeit reeller Funktionen in einer Stelle}

\begin{tcolorbox}[colback=myblue]
\textbf{Definition:}\\
Sei $f\colon \mathbb{R} \to \mathbb{R}$ eine Funktion und $x_0 \in D$. 
Die Funktion $f$ heißt stetig an der Stelle $x_0$, wenn
$$\limeslinks{x}{x_0}f(x)=\limesrechts{x}{x_0}f(x)=f(x_0)$$
gilt.
\end{tcolorbox}

\begin{tcolorbox}[colback=mygreen]
\textbf{Beispiel 2:}\\
Gegeben sei die Funktion
\[
f(x)=x^2
\]

Untersuche die Stetigkeit an der Stelle $x_0=2$.
\end{tcolorbox}

\textbf{Funktionswert in $x_0=2$:}

$$f(2)=2^2=4$$

\textbf{Linker Grenzwert:}

\[
\begin{array}{c|c}
x & f(x)=x^2 \\
\hline
1{,}9 & 3{,}61 \\
1{,}99 & 3{,}9601 \\
1{,}999 & 3{,}996001
\end{array}
\]
Es ist zu vermuten, dass
\[
\limeslinks{x}{2} f(x)=4.
\]

\textbf{Rechter Grenzwert:}

\[
\begin{array}{c|c}
x & f(x)=x^2 \\
\hline
2{,}1 & 4{,}41 \\
2{,}01 & 4{,}0401 \\
2{,}001 & 4{,}004001
\end{array}
\]
Es ist zu vermuten, dass
\[
\limesrechts{x}{2} f(x)=4.
\]


\textbf{Fazit:}

Da

\[
\limeslinks{x}{2}f(x)=\limesrechts{x}{2}f(x)=f(2)=4
\]

gilt, ist die Funktion bei $x=2$ \textbf{stetig}.

\newpage
\begin{tcolorbox}[colback=mygreen]
\textbf{Beispiel 2:}\\
Gegeben sei die Funktion
\[
f(x)=
\begin{cases}
2x+1, & x<1 \\
5, & x=1 \\
x+3, & x>1
\end{cases}
\]

Untersuche die Stetigkeit an der Stelle $x_0=1$.
\end{tcolorbox}

\textbf{Funktionswert in $x_0=1$:}

$$f(1)=5$$

\textbf{Linker Grenzwert}

Wir nähern uns $1$ von links:

\[
\begin{array}{c|c}
x & f(x)=2x+1 \\
\hline
0,9 & 2,8 \\
0,99 & 2,98 \\
0,999 & 2,998
\end{array}
\]
Es ist zu vermuten, dass
\[
\limeslinks{x}{1} f(x)=3.
\]


\textbf{Fazit:}
\vspace{0.5cm}

Da
\[
\limeslinks{x}{1} f(x) \ne f(1),
\]

ist die Funktion bei $x=1$ \textbf{nicht stetig}.
\newpage
\subsection*{Aufgaben}


1. Untersuchen Sie, ob die Funktion $f$ an der Stelle $x_0$ stetig ist.

\begin{enumerate}[label=(\alph*)]

\item $f(x)=4x+3$, \quad $x_0=0$

\item $f(x)=\sqrt{x^2}$, \quad $x_0=0$

\item $f(x)=\dfrac{1}{x}$, \quad $x_0=-1$

\item $f(x)=\dfrac{x-1}{x}$, \quad $x_0=0,01$

\item 
\[
f(x)=
\begin{cases}
\dfrac{x^2-1}{x-1}, & \text{für } x\ne1\\
2, & \text{sonst}
\end{cases},\quad x_0=1
\]


\item
\[
f(x)=
\begin{cases}
4x-5, & \text{für } x\le3\\
2x+1, & \text{für } x>3
\end{cases},\quad x_0=3
\]


\end{enumerate}

\vspace{1cm}

2. Ermitteln Sie, für welchen Wert von $a$ die Funktion $$f(x)=
\begin{cases}
\dfrac{x^3-1}{x-1}, & \text{für } x\ne1\\
2a+3, & \text{sonst}
\end{cases}
$$ an der Stelle $x_0=1$ stetig ist.



\end{document}

